\section{Diagonal and Block Matrix Families}
\label{sec:diag-block}

Diagonal and block families are the closest finite-dimensional models of
coordinatewise and grouped adaptive preconditioning. Their reachability admits a
clean semidefinite characterization. This connects restricted geometric
complexity to the classical literature on scaling and equilibration
\cite{bauer1963optimally,vandersluis1969equilibration,
shapiro1982optimally,shapiro1985block,sinkhorn1967doubly,
knight2013balancing,qu2024optimal}, to condition-number optimization through
convex programming \cite{marechal2009optimizing,lu2011minimizing}, and to
semidefinite feasibility \cite{vandenberghe1996semidefinite,boyd2004convex}.
The role of this section is to place those reachability questions inside the
structured affine-invariant target-distance benchmark and to record primal and
dual certificates in the notation used by \(D_{K,\F}\).

\begin{center}
\begin{tabular}{p{0.28\linewidth}p{0.30\linewidth}p{0.32\linewidth}}
\toprule
Topic & Closest classical object & Role here \\
\midrule
Diagonal/block scaling &
Optimal scaling and equilibration criteria &
Endpoint reachability for a restricted metric family \\
Condition-number minimization &
Variational and convex-programming formulations &
The same target embedded in an affine-invariant path-distance benchmark \\
Semidefinite feasibility &
Primal feasibility and dual separation &
Certificate language for finite \(D_{K,\F}\) or unreachable targets \\
\bottomrule
\end{tabular}
\end{center}

\subsection{Diagonal reachability}

Let
\[
  \F_{\diag}=\{D=\diag(e^{x_1},\ldots,e^{x_d}):x\in\R^d\}.
\]
Inside this family the affine-invariant line element is Euclidean in log
coordinates:
\[
  \Len_{\diag}(x_{[0,T]})=\int_0^T\norm{\dot x_t}_2\,\dd t.
\]
The diagonal target set is
\[
  \mathcal X_K(H)=
  \{x\in\R^d:\kappa(\diag(e^{-x})H)\le K\}.
\]

\begin{lemma}[Diagonal distance identity]
\label{lem:diag-distance-identity}
The map
\[
  x\mapsto H^{-1/2}\diag(e^x)H^{-1/2}
\]
is an isometric realization of the diagonal metric family with its
affine-invariant induced length.  With
\(S_0=H^{-1/2}\diag(e^{x_0})H^{-1/2}\), write
\(D_{K,\diag}(x_0;H):=D_{K,\diag}(S_0;H)\).  The set
\(\mathcal X_K(H)\) is closed, and
\[
  D_{K,\diag}(x_0;H)=\dist_{\ell^2}(x_0,\mathcal X_K(H)).
\]
\end{lemma}

\begin{theorem}[Diagonal LMI reachability]
\label{thm:diag-lmi}
For \(K\ge1\), the following are equivalent:
\begin{enumerate}[leftmargin=*,itemsep=0.1em]
  \item there exists a positive diagonal \(D\) with
  \(\kappa(D^{-1}H)\le K\);
  \item there exists a positive diagonal \(E\) with
  \[
    E\preceq H\preceq K E.
  \]
\end{enumerate}
Consequently,
\[
  K_{\diag}^\ast(H)
  =
  \inf\{K\ge1:\exists E\succ0\ \mathrm{diagonal},\ E\preceq H\preceq KE\}.
\]
For fixed \(K\), diagonal reachability is an LMI feasibility problem.
\end{theorem}

\begin{remark}[Certificate interpretation]
\Cref{thm:diag-lmi} separates reachability from path length. If the LMI is
infeasible, then \(D_{K,\diag}(x_0;H)=+\infty\) for every initial diagonal
metric. If it is feasible, the remaining problem is the intrinsic distance from
the initial log-scale \(x_0\) to the feasible set \(\mathcal X_K(H)\).
Semidefinite duality can therefore provide certificates that no diagonal
preconditioner can reach a desired condition-number threshold.
\end{remark}

\begin{corollary}[Aligned diagonal optimum]
\label{cor:diag-aligned}
If \(H=\diag(h_1,\ldots,h_d)\) and \(G_0\) is diagonal, then
\[
  D_{K,\diag}(S_0;H)=D_K(S_0).
\]
\end{corollary}

Thus diagonal geometry can be optimal when the target curvature is already
aligned with the coordinate axes. Its loss comes from directional mismatch.

\subsection{Block reachability}

Let the coordinates be split as
\[
  \R^d=V_1\oplus\cdots\oplus V_m,
\qquad \dim V_j=d_j,
\]
and let \(\F_{\mathrm{block}}\) be the corresponding block-diagonal positive
definite family.

\begin{theorem}[Block LMI reachability]
\label{thm:block-lmi}
For \(K\ge1\), the following are equivalent:
\begin{enumerate}[leftmargin=*,itemsep=0.1em]
  \item there exists \(G\in\F_{\mathrm{block}}\) with
  \(\kappa(G^{-1}H)\le K\);
  \item there exists \(E\in\F_{\mathrm{block}}\) with
  \[
    E\preceq H\preceq K E.
  \]
\end{enumerate}
Therefore block reachability is also an LMI feasibility problem.
\end{theorem}

\subsection{Dual infeasibility certificates}

The LMI formulation also gives a checkable certificate when a target condition
number is impossible. Let \(\mathcal L\) be either the diagonal symmetric
subspace or a fixed block-diagonal symmetric subspace, and let
\(\Pi_{\mathcal L}\) be the Frobenius-orthogonal projection onto
\(\mathcal L\). For fixed \(K\), the primal reachability problem asks for
\(E\in\mathcal L\) such that
\[
  H-E\succeq0,\qquad KE-H\succeq0.
\]

\begin{proposition}[Structured SDP dual certificate]
\label{prop:structured-dual-certificate}
If there exist \(P,Q\in\PSD^d\) such that
\[
  \Pi_{\mathcal L}(P-KQ)=0,
  \qquad
  \ip{P-Q}{H}<0,
\]
then no \(E\in\mathcal L\) satisfies \(E\preceq H\preceq KE\). Hence the
diagonal or block restricted target is empty at threshold \(K\).
\end{proposition}

\begin{corollary}[Sign-flip rank-one certificate]
\label{cor:sign-flip-certificate}
Let
\[
  T=\blockdiag(s_1 I_{d_1},\ldots,s_m I_{d_m}),
  \qquad s_j\in\{-1,1\},
\]
and let \(\mathcal L\) be the corresponding block-diagonal subspace. If
\[
  \lambda_{\min}(K THT-H)<0,
\]
then the block family cannot reach condition number \(K\). The diagonal case is
the special case \(d_j=1\).
\end{corollary}

The sign-flip certificate is not meant to solve every infeasible SDP. It gives
a cheap, inspectable dual witness for synthetic examples where the obstruction
is cross-block or cross-coordinate coupling.

\begin{proposition}[Block Jacobi upper bound]
\label{prop:block-jacobi}
Let
\[
  G_{\mathrm{BJ}}=\blockdiag(H_{11},\ldots,H_{mm})
\]
with every \(H_{jj}\succ0\), and set
\[
  \widetilde H=G_{\mathrm{BJ}}^{-1/2}HG_{\mathrm{BJ}}^{-1/2}.
\]
Then
\[
  K_{\mathrm{block}}^\ast(H)\le \kappa(\widetilde H).
\]
If \(\opnorm{\widetilde H-I}\le\delta<1\), then
\[
  K_{\mathrm{block}}^\ast(H)\le\frac{1+\delta}{1-\delta}.
\]
\end{proposition}

The block family can absorb all within-block curvature. The residual
difficulty is the normalized off-block interaction.
